\documentclass[11pt]{article}
\usepackage{booktabs}
\usepackage[margin=1in]{geometry}
\usepackage{fontspec}
\setmainfont{lmroman10-regular.otf}[BoldFont=lmroman10-bold.otf,ItalicFont=lmroman10-italic.otf,BoldItalicFont=lmroman10-bolditalic.otf]
\setsansfont{lmsans10-regular.otf}[BoldFont=lmsans10-bold.otf,ItalicFont=lmsans10-oblique.otf]
\setmonofont{lmmono10-regular.otf}[ItalicFont=lmmono10-italic.otf]
\title{Lua Lm Fontspec}
\author{A. Author}
\date{1 January 2026}
\begin{document}
\maketitle
\begin{abstract}
This synthetic benchmark document exercises the engine with typical content.
\end{abstract}
\tableofcontents
\section{Average Impact}\label{sec:1}

We find that the modest uncertainty relates the parameter, while the final value affects the general pattern. The empirical system limits the observed sector of the benefit. The implicit schedule improves the optimal structure of the decision (see Section~\ref{sec:5}). We find that the modest measure stabilizes the size, while the complex price conditions the observed correlation (see Section~\ref{sec:6}). We find that the active finding bounds the pressure, while the low impact limits the observed selection. A low pressure shapes each value in the implicit cluster (see Section~\ref{sec:4}).

The broad variable shapes the structural distribution of the context. We find that the persistent schedule varies the variance, while the possible observation affects the optimal stability. A local price relates each control in the dynamic assumption. A large strategy stabilizes each limit in the final context. We find that the direct control determines the hypothesis, while the common response drives the implicit condition.

\section{Average Information}\label{sec:2}

In the early shock, the approach stabilizes a average growth and the size (see Section~\ref{sec:8}). We find that the random coefficient stabilizes the choice, while the high profit explains the narrow return. We find that the typical rate changes the market, while the common rate improves the general process. In the persistent supply, the margin limits a direct constraint and the result. The random condition changes the marginal weight of the solution. The observed result relates the modest channel of the forecast.

In the stable assumption, the index conditions a robust measure and the test. The low network reduces the typical utility of the behavior. In the early correlation, the risk supports a dynamic gain and the test. A robust structure conditions each latency in the average regression. The observed framework reduces the low channel of the market.

\section{Final Assumption}\label{sec:3}

A complex curve affects each sample in the low layer. The partial policy conditions the efficient variance of the state. We find that the clear interaction reflects the regime, while the stable trade stabilizes the robust rate. A broad risk explains each test in the low limit. A broad trade stabilizes each size in the direct pattern. We find that the global condition explains the information, while the common framework determines the typical regime (see Section~\ref{sec:5}).

\begin{table}[tbp]\centering\caption{Local curve summary.}\label{tab:b3}\begin{tabular}{lrrr}\toprule Gain & Uncertainty & Efficiency & Interval \\ \midrule
experiment & 90.97 & 97.71 & 73.61 \\
comparison & 82.56 & 71.52 & 37.95 \\
impact & 26.03 & 96.25 & 74.56 \\
distribution & 58.30 & 58.90 & 51.95 \\
scale & 38.54 & 70.11 & 67.04 \\
selection & 29.81 & 22.75 & 73.63 \\
knowledge & 14.66 & 34.72 & 31.40 \\
stability & 69.26 & 96.20 & 93.55 \\
\bottomrule\end{tabular}\end{table}

Table~\ref{tab:b3} lists the values. The general feature reflects the strong labor of the constraint. We find that the active distribution changes the investment, while the large measure tracks the partial uncertainty.

We find that the central pressure limits the probability, while the implicit method bounds the random efficiency. The sufficient index determines the simple model of the state. We find that the observed condition bounds the performance, while the empirical system supports the marginal scale. In the general approach, the latency varies a primary portfolio and the delay. A strong analysis changes each risk in the broad risk.

\section{Marginal Layer}\label{sec:4}

A joint assumption supports each judgment in the global cluster. We find that the structural decision stabilizes the threshold, while the possible scale tracks the low policy. In the early information, the selection bounds a persistent error and the decision. The marginal gain explains the common margin of the noise. We find that the complex trend relates the source, while the local market reduces the optimal process. A structural process changes each impact in the primary variance.

In the modest factor, the risk conditions a large latency and the outcome\footnote{A stable return reduces each constraint in the sharp population. We find that the dynamic coefficient supports the size, while the sharp analysis drives the primary flow.}. In the persistent trade, the source changes a central impact and the assumption\footnote{The partial curve reflects the central regime of the interaction. In the high knowledge, the population shapes a random knowledge and the knowledge.}. In the efficient gain, the design affects a partial equilibrium and the policy\footnote{We find that the average labor improves the test, while the sharp stability affects the general performance. The local curve reduces the implicit balance of the index.}. In the final feature, the effect explains a common probability and the volume\footnote{The clear state varies the implicit result of the trade. A central loss conditions each impact in the sharp efficiency.}. We find that the narrow demand determines the efficiency, while the persistent forecast limits the random production\footnote{A typical approach changes each threshold in the stable equilibrium. We find that the general balance changes the market, while the modest loss reflects the general regression.}. The implicit analysis reduces the clear context of the production\footnote{The broad model tracks the direct weight of the prediction. In the joint shock, the gain stabilizes a dynamic regime and the scale.}.

The benchmark edit marker reads BENCH-A.

In the partial risk, the state reflects a active experiment and the strategy. A dynamic distribution changes each response in the direct control. In the broad utility, the response affects a general response and the schedule. In the clear risk, the performance determines a high analysis and the size. A direct feature affects each pressure in the low information.

\section{Strong Judgment}\label{sec:5}

A persistent result reflects each price in the typical response. A simple response varies each shock in the simple growth (see Section~\ref{sec:8}). The implicit solution relates the final source of the comparison. A low weight predicts each technique in the direct regime. The active state drives the modest feature of the schedule (see Section~\ref{sec:8}). The active effect explains the direct spread of the constraint.

In the empirical assumption, the selection varies a direct error and the rate. A empirical spread conditions each labor in the common schedule. In the narrow supply, the efficiency drives a sharp efficiency and the channel. In the general framework, the context drives a dynamic test and the trend. The stable information relates the high cluster of the portfolio.

\section{Possible Network}\label{sec:6}

We find that the local trade determines the hypothesis, while the partial approach explains the joint context. We find that the dynamic risk drives the factor, while the local growth supports the central finding. We find that the constant behavior tracks the outcome, while the persistent shock shapes the complex knowledge. In the dynamic value, the context affects a observed coefficient and the curve. A high limit conditions each regression in the dynamic comparison. A broad prediction affects each scale in the broad spread.

\begin{table}[tbp]\centering\caption{Early coefficient summary.}\label{tab:b6}\begin{tabular}{lrrr}\toprule Equilibrium & Test & Trend & Variable \\ \midrule
policy & 47.91 & 88.22 & 89.41 \\
efficiency & 50.24 & 89.47 & 27.11 \\
comparison & 30.04 & 97.30 & 5.92 \\
trend & 6.25 & 72.39 & 7.30 \\
behavior & 5.69 & 23.00 & 0.66 \\
strategy & 38.33 & 51.49 & 98.62 \\
curve & 82.35 & 23.30 & 86.91 \\
channel & 3.56 & 82.29 & 84.04 \\
\bottomrule\end{tabular}\end{table}

Table~\ref{tab:b6} lists the values. In the empirical price, the distribution determines a empirical knowledge and the assumption. In the general weight, the capital determines a early correlation and the control.

In the average trend, the finding limits a local distribution and the coefficient. A complex loss relates each demand in the robust trade. The possible utility relates the modest regression of the regime. The marginal approach bounds the common distribution of the judgment. We find that the random function limits the value, while the constant range changes the persistent sector.

\section{Structural Variance}\label{sec:7}

In the direct balance, the supply changes a constant technique and the parameter (see Section~\ref{sec:1}). We find that the primary dynamic improves the limit, while the empirical factor varies the large cost. In the sharp sector, the constraint varies a constant rate and the price. A common data varies each loss in the sharp cost (see Section~\ref{sec:2}). We find that the simple income changes the profit, while the complex sector reflects the sufficient strategy. A robust sector explains each error in the simple parameter (see Section~\ref{sec:1}).

A constant cluster reduces each pattern in the marginal price. We find that the possible information explains the flow, while the joint pattern predicts the partial signal. In the high stability, the benefit drives a common performance and the growth. We find that the observed selection improves the labor, while the structural result varies the early sample. A modest cluster bounds each process in the direct gain.

\section{Possible Example}\label{sec:8}

A constant interaction stabilizes each layer in the strong correlation. The general result changes the early judgment of the decision. A average trend bounds each probability in the random correlation. We find that the sufficient shock changes the weight, while the dynamic behavior conditions the direct data. We find that the partial variable tracks the selection, while the optimal approach relates the random sector. We find that the sufficient behavior reflects the index, while the possible feature supports the stable investment.

A common signal reflects each analysis in the primary efficiency\footnote{We find that the simple channel reflects the parameter, while the global impact bounds the central uncertainty. A efficient income drives each distribution in the possible noise.}. We find that the final policy tracks the technique, while the final judgment reflects the complex demand\footnote{In the clear capital, the pattern drives a sharp spread and the condition. A sufficient market affects each range in the direct threshold.}. A simple impact conditions each loss in the primary quality\footnote{We find that the strong production explains the regression, while the general technique limits the observed loss. A empirical model improves each cluster in the local system.}. A complex technique determines each production in the large model\footnote{The constant impact relates the partial structure of the variance. A partial prediction supports each channel in the sufficient state.}. The random constraint shapes the primary scale of the uncertainty\footnote{We find that the empirical regime limits the behavior, while the low variable predicts the local supply. A optimal variable tracks each price in the stable population.}. In the active return, the shock varies a local behavior and the network\footnote{A robust income improves each value in the sharp growth. A complex income stabilizes each feature in the marginal finding.}.

The early test shapes the strong data of the feature. A efficient process improves each balance in the random process. The structural analysis predicts the sufficient performance of the prediction. A clear analysis improves each income in the common delay. We find that the stable solution varies the trend, while the partial comparison conditions the modest forecast.

\end{document}
